\documentclass[a4paper,11pt]{ctexart}
\usepackage{amsmath, amssymb, amsfonts}
\usepackage{geometry}
\geometry{left=1.5cm, right=1.5cm, top=2.0cm, bottom=2.0cm, headsep=0.3cm, footskip=1cm}
\usepackage{listings}
\usepackage{xcolor}
\usepackage{tikz}
\usepackage{pgfplots}
\usepackage{hyperref}
\usepackage{fancyhdr}
\usepackage{tcolorbox}
\tcbuselibrary{breakable, skins}
\usepackage{array}
\usepackage{booktabs}
\usepackage[shortlabels]{enumitem}
\usepackage{needspace}
\usepackage{multirow}

\AtBeginDocument{\hypersetup{colorlinks=true, linkcolor=blue, filecolor=magenta, urlcolor=cyan}}
\setlist{nosep, itemsep=0.8ex, parsep=0.1em}
\setlist[itemize]{leftmargin=1.8em}
\setlist[enumerate]{leftmargin=1.8em}

\definecolor{codegreen}{rgb}{0,0.6,0}
\definecolor{codegray}{rgb}{0.5,0.5,0.5}
\definecolor{codepurple}{rgb}{0.58,0,0.82}
\definecolor{codebg}{rgb}{0.98,0.98,0.98}
\definecolor{tipsblue}{rgb}{0.85,0.93,1.0}
\definecolor{highlightyellow}{rgb}{1.0,0.97,0.8}

\lstdefinestyle{mystyle}{backgroundcolor=\color{codebg},commentstyle=\color{codegreen},keywordstyle=\color{magenta}\bfseries,numberstyle=\tiny\color{codegray},stringstyle=\color{codepurple},basicstyle=\ttfamily\small,breakatwhitespace=false,breaklines=true,captionpos=b,keepspaces=true,numbers=left,numbersep=6pt,showspaces=false,showstringspaces=false,showtabs=false,tabsize=2,language=Python,frame=single,framerule=0.5pt,rulecolor=\color{gray!40}}
\lstset{style=mystyle}

\newtcolorbox{problembox}[1][]{colback=white,colframe=blue!60!black,title=\textbf{#1},fonttitle=\bfseries\small,boxrule=1pt,arc=3pt,outer arc=3pt,coltitle=black,before skip=10pt,after skip=8pt}
\newtcolorbox{solutionbox}[1][]{enhanced,breakable,colback=green!3!white,colframe=green!50!black,title=\textbf{#1},fonttitle=\bfseries\small,boxrule=1pt,arc=3pt,outer arc=3pt,coltitle=black,before skip=8pt,after skip=10pt}
\newtcolorbox{metaphorbox}[1][]{colback=orange!5!white,colframe=orange!70!black,title=\textbf{#1},fonttitle=\bfseries\small,boxrule=1pt,arc=3pt,outer arc=3pt,coltitle=black,before skip=8pt,after skip=8pt}
\newtcolorbox{beginnerbox}[1][]{colback=tipsblue,colframe=blue!50!black,title=\textbf{#1},fonttitle=\bfseries\small,boxrule=1pt,arc=3pt,outer arc=3pt,coltitle=black,before skip=8pt,after skip=8pt}
\newtcolorbox{appbox}[1][]{colback=green!3!white,colframe=green!60!black,title=\textbf{#1},fonttitle=\bfseries\small,boxrule=1pt,arc=3pt,outer arc=3pt,coltitle=black,before skip=8pt,after skip=10pt}
\newtcolorbox{keybox}[1][]{colback=highlightyellow,colframe=orange!80!black,title=\textbf{#1},fonttitle=\bfseries\small,boxrule=1pt,arc=3pt,outer arc=3pt,coltitle=black,before skip=8pt,after skip=8pt}

\pagestyle{fancy}
\fancyhf{}
\fancyhead[R]{机器学习-第5章}
\fancyhead[L]{KNN算法 · 练习题}
\fancyfoot[C]{\thepage}
\addtolength{\headheight}{2pt}

\parskip=2pt plus 1pt minus 0.5pt
\linespread{1.35}
\raggedbottom
\widowpenalty=10000
\clubpenalty=10000

\title{\textbf{机器学习\\第5章\quad KNN算法}\\[0.5em]
{\large\color{gray}——最简单的算法：近朱者赤，近墨者黑}}
\date{\today}

\begin{document}
\maketitle

\begin{beginnerbox}[本章题目导览：你将挑战哪些核心问题？]
    KNN（K-Nearest Neighbors）是机器学习中最直观、最易理解的算法之一——"你的 $k$ 个最近邻是什么类别，你就是什么类别"。本章 5 道题覆盖 KNN 的完整体系：
    \begin{itemize}
        \item \textbf{Q5-01（原理推导）}：KNN 工作原理手动计算——给定数据点，算距离，找邻居，投票决策。
        \item \textbf{Q5-02（参数分析）}：$k$ 值如何影响决策边界？过大/过小的 $k$ 分别导致什么问题？
        \item \textbf{Q5-03（距离度量）}：欧氏距离、曼哈顿距离、闵可夫斯基距离的计算与比较。
        \item \textbf{Q5-04（API 实践）}：sklearn KNN 分类与回归的参数解析与代码实现。
        \item \textbf{Q5-05（案例综合）}：心脏病预测——完整的 KNN 建模流程，包括特征工程、超参数调优。
    \end{itemize}
\end{beginnerbox}

\section{第 5 章\quad KNN：最直观的机器学习算法}

\begin{metaphorbox}[先建立一个总感觉：KNN 在做什么？]
    KNN 的逻辑极其简单，就是现实生活中的一句话：\textbf{"近朱者赤，近墨者黑。"}\\
    如果你不知道一个人是好人还是坏人，就看他最亲近的 $k$ 个朋友——如果大部分朋友是好人，那他大概率也是好人。\\
    KNN 没有显式的"训练过程"：它把整个训练集"记住"，预测时直接找最近邻来投票。这叫做\textbf{懒惰学习（Lazy Learning）}。
\end{metaphorbox}

%% ============================================================
\Needspace{5\baselineskip}
\begin{problembox}[Q5-01\quad KNN 工作原理：手动计算分类]
    \textbf{题目描述：}\\
    训练集有以下 6 个数据点（2 个特征），标签为 A 类或 B 类：

    \begin{center}
    \begin{tabular}{cccc}
        \toprule
        编号 & $x_1$ & $x_2$ & 类别 \\
        \midrule
        1 & 1 & 1 & A \\
        2 & 2 & 2 & A \\
        3 & 3 & 1 & A \\
        4 & 6 & 5 & B \\
        5 & 7 & 4 & B \\
        6 & 5 & 6 & B \\
        \bottomrule
    \end{tabular}
    \end{center}

    新样本为 $Q = (4, 3)$，使用 KNN 算法（欧氏距离）进行分类。

    \vspace{0.2cm}
    \textbf{问题：}
    \begin{enumerate}[(1)]
        \item 计算 $Q$ 到所有 6 个训练点的欧氏距离（保留两位小数），填写下表：

        \begin{center}
        \begin{tabular}{cccc}
            \toprule
            训练点 & 坐标 & 到 $Q(4,3)$ 的距离 & 类别 \\
            \midrule
            1 & (1,1) & \underline{\quad} & A \\
            2 & (2,2) & \underline{\quad} & A \\
            3 & (3,1) & \underline{\quad} & A \\
            4 & (6,5) & \underline{\quad} & B \\
            5 & (7,4) & \underline{\quad} & B \\
            6 & (5,6) & \underline{\quad} & B \\
            \bottomrule
        \end{tabular}
        \end{center}

        \item 当 $k=3$ 时，找出 3 个最近邻，并通过投票确定 $Q$ 的预测类别。
        \item 当 $k=5$ 时，预测类别是否改变？
        \item KNN 用于\textbf{回归}时，预测方式与分类有何不同？设 6 个点的目标值分别为 $[10, 12, 11, 30, 28, 31]$，$k=3$ 时 $Q$ 的预测值是多少？
    \end{enumerate}
    {\color{gray}\small\textit{来源溯源：[文档第 5 章 5.1 KNN算法介绍]}}
\end{problembox}

\begin{solutionbox}[参考答案与详细解析]
    \textbf{(1) 距离计算（$d = \sqrt{(x_1-4)^2+(x_2-3)^2}$）：}

    \begin{center}
    \begin{tabular}{cccc}
        \toprule
        训练点 & 坐标 & 距离 & 类别 \\
        \midrule
        1 & (1,1) & $\sqrt{9+4}=\sqrt{13}\approx3.61$ & A \\
        2 & (2,2) & $\sqrt{4+1}=\sqrt{5}\approx2.24$ & A \\
        3 & (3,1) & $\sqrt{1+4}=\sqrt{5}\approx2.24$ & A \\
        4 & (6,5) & $\sqrt{4+4}=\sqrt{8}\approx2.83$ & B \\
        5 & (7,4) & $\sqrt{9+1}=\sqrt{10}\approx3.16$ & B \\
        6 & (5,6) & $\sqrt{1+9}=\sqrt{10}\approx3.16$ & B \\
        \bottomrule
    \end{tabular}
    \end{center}

    \textbf{(2) $k=3$：}

    最近 3 个点（按距离排序）：点2（2.24，A），点3（2.24，A），点4（2.83，B）。投票：A 得 2 票，B 得 1 票。\textbf{预测为 A 类}。

    \textbf{(3) $k=5$：}

    最近 5 个点：点2（A），点3（A），点4（B），点5（B），点6（B）。投票：A 得 2 票，B 得 3 票。\textbf{预测变为 B 类}——$k$ 的选择对结果有重要影响。

    \textbf{(4) KNN 回归：}

    KNN 回归预测新样本的输出为 $k$ 个最近邻目标值的\textbf{平均值}（而非投票）。$k=3$ 时最近 3 个点为点2（目标值12）、点3（11）、点4（30）：
    \[
    \hat{y} = \frac{12+11+30}{3} = \frac{53}{3} \approx \mathbf{17.67}
    \]
\end{solutionbox}

%% ============================================================
\Needspace{5\baselineskip}
\begin{problembox}[Q5-02\quad $k$ 值的影响：过大/过小的后果]
    \textbf{题目描述：}\\
    KNN 中最关键的超参数是 $k$（近邻数），它直接控制模型复杂度。

    \vspace{0.2cm}
    \textbf{问题：}
    \begin{enumerate}[(1)]
        \item 当 $k=1$ 时，KNN 的决策边界是什么形状？它的训练误差是多少？为什么？

        \item 当 $k=n$（$n$ 为训练集样本总数）时，KNN 的预测结果是什么？这对应欠拟合还是过拟合？

        \item 随着 $k$ 从 1 增大到 $n$，模型的偏差和方差分别如何变化？（描述趋势）

        \item 在实际使用中，通常如何选择最优的 $k$？写出用 sklearn 的 \texttt{GridSearchCV} 进行 $k$ 值网格搜索的核心代码。

        \item KNN 的时间复杂度是多少（预测一个新样本时）？当训练集很大时，有什么加速策略？
    \end{enumerate}
    {\color{gray}\small\textit{来源溯源：[文档第 5 章 5.1 关键参数 / 5.2 API使用]}}
\end{problembox}

\begin{solutionbox}[参考答案与详细解析]
    \textbf{(1) $k=1$ 时：}

    决策边界是 Voronoi 图（每个训练点控制其周围最近的区域），边界极度不规则（锯齿状）。训练误差\textbf{精确为 0}——因为每个训练点的最近邻就是它自己，预测结果与真实标签完全一致。这是典型的过拟合：决策边界高度拟合训练数据的噪声，泛化能力差。

    \textbf{(2) $k=n$ 时：}

    所有训练样本都是近邻，预测结果是训练集中\textbf{多数类别}（对分类）或\textbf{所有训练样本目标值的均值}（对回归）——对所有新样本一视同仁，完全忽略输入特征，\textbf{对应欠拟合}（高偏差，低方差）。

    \textbf{(3) 偏差-方差随 $k$ 变化：}

    $k$ 从 1 增大到 $n$：偏差\textbf{单调增大}（平滑化导致决策边界越来越粗略）；方差\textbf{单调减小}（平均更多邻居，对单个噪声点不敏感）。最优 $k$ 在偏差-方差权衡的"最低总误差"处。

    \textbf{(4) 网格搜索选 $k$：}
    \begin{lstlisting}[language=Python]
from sklearn.neighbors import KNeighborsClassifier
from sklearn.model_selection import GridSearchCV

param_grid = {'n_neighbors': range(1, 31, 2)}  # 搜索奇数k，避免平票
knn = KNeighborsClassifier()
grid_search = GridSearchCV(knn, param_grid, cv=5, scoring='accuracy')
grid_search.fit(X_train_scaled, y_train)
print(f"最优k: {grid_search.best_params_}")
print(f"最优CV准确率: {grid_search.best_score_:.3f}")
    \end{lstlisting}

    \textbf{(5) 时间复杂度与加速：}

    朴素 KNN 预测一个新样本需要计算它与所有 $n$ 个训练样本的距离，复杂度为 $O(nd)$（$d$ 为维度），总预测复杂度 $O(nmd)$（$m$ 个测试样本）。训练集大时极慢。

    \begin{keybox}[KNN 加速数据结构]
        \textbf{KD-Tree（K-Dimensional Tree）}：将训练数据组织为二叉树，预测时只搜索邻近分支，平均复杂度 $O(d\log n)$，但高维时退化。\\
        \textbf{Ball-Tree}：用超球体组织数据，对高维数据比 KD-Tree 更鲁棒。\\
        sklearn 的 \texttt{KNeighborsClassifier(algorithm='auto')} 会自动选择合适的数据结构。
    \end{keybox}
\end{solutionbox}

%% ============================================================
\Needspace{5\baselineskip}
\begin{problembox}[Q5-03\quad 距离度量：欧氏、曼哈顿、闵可夫斯基]
    \textbf{题目描述：}\\
    KNN 的核心操作是计算距离，距离度量的选择影响"近邻"的定义，从而影响分类结果。

    \vspace{0.2cm}
    \textbf{题目数据：}两个点 $P=(1,2,3)$，$Q=(4,0,6)$（三维空间）。

    \vspace{0.2cm}
    \textbf{问题：}
    \begin{enumerate}[(1)]
        \item 计算 $P$ 和 $Q$ 之间的：
        \begin{enumerate}[a.]
            \item \textbf{欧氏距离}（$L_2$ 范数）：$d = \sqrt{\sum_i(p_i-q_i)^2}$
            \item \textbf{曼哈顿距离}（$L_1$ 范数）：$d = \sum_i|p_i-q_i|$
            \item \textbf{切比雪夫距离}（$L_\infty$ 范数）：$d = \max_i|p_i-q_i|$
        \end{enumerate}

        \item 闵可夫斯基距离（Minkowski Distance）的公式是 $d = \left(\sum_i|p_i-q_i|^r\right)^{1/r}$，当 $r=1, 2, \infty$ 时分别对应上面哪种距离？

        \item 填写下表，对比三种距离的特点与适用场景：
        \begin{center}
        \begin{tabular}{p{2.5cm} p{4cm} p{5cm}}
            \toprule
            距离度量 & 几何直觉 & 适用场景 \\
            \midrule
            欧氏距离 & \underline{\hspace{3.5cm}} & 特征无明显方向偏好；连续数值特征 \\[1.2em]
            曼哈顿距离 & \underline{\hspace{3.5cm}} & \underline{\hspace{4.5cm}} \\[1.2em]
            切比雪夫距离 & \underline{\hspace{3.5cm}} & \underline{\hspace{4.5cm}} \\
            \bottomrule
        \end{tabular}
        \end{center}

        \item 在 KNN 中使用距离度量时，为什么\textbf{必须先对特征进行归一化或标准化}？用具体数字说明不标准化的危害。
    \end{enumerate}
    {\color{gray}\small\textit{来源溯源：[文档第 5 章 5.3 常见距离度量方法]}}
\end{problembox}

\begin{solutionbox}[参考答案与详细解析]
    \textbf{(1) 距离计算（$P=(1,2,3), Q=(4,0,6)$，差值为 $[3, -2, 3]$）：}
    \begin{enumerate}[a.]
        \item \textbf{欧氏距离}：$\sqrt{3^2+(-2)^2+3^2} = \sqrt{9+4+9} = \sqrt{22} \approx \mathbf{4.69}$
        \item \textbf{曼哈顿距离}：$|3|+|-2|+|3| = 3+2+3 = \mathbf{8}$
        \item \textbf{切比雪夫距离}：$\max(3, 2, 3) = \mathbf{3}$
    \end{enumerate}

    \textbf{(2) 闵可夫斯基的统一框架：}

    $r=1$：曼哈顿距离；$r=2$：欧氏距离；$r\to\infty$：切比雪夫距离（极限情况）。

    \textbf{(3) 特点与适用场景：}
    \begin{center}
    \begin{tabular}{p{2.5cm} p{4cm} p{5cm}}
        \toprule
        距离度量 & 几何直觉 & 适用场景 \\
        \midrule
        欧氏距离 & 两点之间的直线距离（空中飞行） & 连续特征，无明显方向性；图像像素距离 \\[0.8em]
        曼哈顿距离 & 只能沿轴走的出租车距离（城市街道） & 特征独立性强，稀疏高维数据（如文本 TF-IDF） \\[0.8em]
        切比雪夫距离 & 棋盘上国王走一步能到的范围 & 当最大单维差异是主要关注点时；棋类 AI \\
        \bottomrule
    \end{tabular}
    \end{center}

    \textbf{(4) 不标准化的危害（具体数字示例）：}

    假设特征1是年龄（20--60，范围 40），特征2是年收入（20000--200000，范围 180000）。不标准化时，两点 $A=(30, 50000)$ 和 $B=(50, 50000)$ 的欧氏距离 $= \sqrt{400 + 0} = 20$；而 $A=(30, 50000)$ 和 $C=(30, 51000)$ 的距离 $= \sqrt{0 + 1000000} = 1000$。年收入变化 1000 元的距离竟是年龄变化 20 岁的 50 倍——年龄特征\textbf{完全被淹没}，KNN 的"近邻"完全由年收入决定，年龄对结果毫无影响，这与实际业务含义严重不符。
\end{solutionbox}

%% ============================================================
\Needspace{5\baselineskip}
\begin{problembox}[Q5-04\quad KNN API 与完整案例：心脏病预测]
    \textbf{题目描述：}\\
    下面是使用 sklearn KNN 进行心脏病预测的部分代码框架，请结合代码回答问题。

    \begin{lstlisting}[language=Python]
import pandas as pd
from sklearn.model_selection import train_test_split, GridSearchCV
from sklearn.preprocessing import StandardScaler
from sklearn.neighbors import KNeighborsClassifier
from sklearn.metrics import classification_report, accuracy_score
import joblib

# 1. 加载数据
df = pd.read_csv('heart_disease.csv')
X = df.drop('target', axis=1)
y = df['target']

# 2. 数据集划分
X_train, X_test, y_train, y_test = train_test_split(
    X, y, test_size=0.2, random_state=42, stratify=y)

# 3. 特征工程
scaler = StandardScaler()
X_train_scaled = scaler.fit_transform(X_train)
X_test_scaled  = scaler.transform(X_test)   # 注意：只transform，不fit

# 4. 模型训练
knn = KNeighborsClassifier(n_neighbors=5, metric='minkowski', p=2)
knn.fit(X_train_scaled, y_train)

# 5. 评估
y_pred = knn.predict(X_test_scaled)
print(classification_report(y_test, y_pred))

# 6. 超参数调优
param_grid = {'n_neighbors': range(1, 21), 'weights': ['uniform', 'distance']}
grid_search = GridSearchCV(KNeighborsClassifier(), param_grid, cv=5)
grid_search.fit(X_train_scaled, y_train)

# 7. 保存模型
joblib.dump(grid_search.best_estimator_, 'knn_model.pkl')
    \end{lstlisting}

    \vspace{0.2cm}
    \textbf{问题：}
    \begin{enumerate}[(1)]
        \item 解释第 3 步中，为什么测试集只能用 \texttt{transform} 而不能用 \texttt{fit\_transform}？

        \item \texttt{KNeighborsClassifier} 中 \texttt{metric='minkowski', p=2} 等价于哪种距离度量？如果改为 \texttt{p=1} 呢？

        \item \texttt{weights='distance'} 和 \texttt{weights='uniform'} 有什么区别？各在什么情况下更好？

        \item 第 6 步的网格搜索中，搜索空间一共有多少种参数组合？使用 5 折交叉验证，总共需要训练多少个模型？

        \item \texttt{joblib.dump} 和 \texttt{joblib.load} 的用途是什么？加载模型后，如何对一个新样本进行预测？（写出关键代码）
    \end{enumerate}
    {\color{gray}\small\textit{来源溯源：[文档第 5 章 5.2 API使用 / 5.5 案例：心脏病预测]}}
\end{problembox}

\begin{solutionbox}[参考答案与详细解析]
    \textbf{(1) 只 transform 测试集的原因：}

    \texttt{fit\_transform} 会\textbf{重新计算}均值和标准差，测试集有不同的统计量，用测试集统计量缩放后的数据与训练集缩放后的数据\textbf{不在同一尺度空间}，KNN 计算的距离会失去意义（就像训练时用公制、预测时用英制）。正确做法是用训练集的统计量（\texttt{scaler} 已记住）对测试集做相同的缩放。

    \textbf{(2) 距离等价：}

    \texttt{metric='minkowski', p=2} 等价于\textbf{欧氏距离}（$L_2$）；\texttt{p=1} 等价于\textbf{曼哈顿距离}（$L_1$）。

    \textbf{(3) weights 参数：}

    \texttt{uniform}（默认）：所有近邻\textbf{等权重}投票，简单民主；\texttt{distance}：近邻权重与距离\textbf{成反比}（越近权重越高），如 $w_i = 1/d_i$。\texttt{distance} 在边界样本较多时更准确（近的邻居更有参考价值）；若数据噪声较多，\texttt{uniform} 更稳健（避免单个非常近的噪声点主导）。

    \textbf{(4) 搜索空间：}

    $n\_neighbors$ 有 20 个取值，$weights$ 有 2 个取值，总共 $20 \times 2 = 40$ 种参数组合。5 折交叉验证每组参数训练 5 个模型，共训练 $40 \times 5 = \mathbf{200}$ 个模型。

    \textbf{(5) 模型保存与预测：}

    \texttt{joblib.dump} 将训练好的模型对象序列化保存到磁盘，\texttt{joblib.load} 加载回来，无需重新训练：
    \begin{lstlisting}[language=Python]
import joblib
import numpy as np

# 加载模型和scaler（两个都要保存！）
model = joblib.load('knn_model.pkl')
scaler = joblib.load('scaler.pkl')  # 记得也要保存scaler

# 预测新样本（必须用同一个scaler处理）
x_new = np.array([[63, 1, 3, 145, 233, 1, 0, 150, 0, 2.3, 0, 0, 1]])
x_new_scaled = scaler.transform(x_new)
pred = model.predict(x_new_scaled)
pred_proba = model.predict_proba(x_new_scaled)
print(f"预测类别: {pred[0]}, 患病概率: {pred_proba[0][1]:.3f}")
    \end{lstlisting}

    \begin{keybox}[工程提醒：scaler 和 model 必须配套保存]
        很多初学者只保存模型忘记保存 scaler——部署时用了不同的缩放参数，预测结果会完全错误。实际项目中，\texttt{scaler}、\texttt{label encoder}、\texttt{model} 等所有预处理对象都要保存。
    \end{keybox}
\end{solutionbox}

\begin{appbox}[现实应用场景]
    \textbf{KNN 的实际使用场景与局限：}\\
    KNN 在小规模数据（几千到几万条）、特征工程充分的场景下表现良好，常用于推荐系统的 item-based 协同过滤、异常检测（距所有训练点都很远的点是异常）、以及作为快速验证基线。其主要局限：预测慢（需遍历训练集）；高维失效（维度诅咒：高维空间中距离失去意义）；内存占用大（需存储全部训练集）。这些局限使其在工业大规模场景中较少单独使用。
\end{appbox}

\section{本章小结}
\begin{itemize}
    \item \textbf{KNN 核心思想}：懒惰学习，预测时找 $k$ 个最近邻投票/取均值——简单但计算代价高。
    \item \textbf{$k$ 值权衡}：$k$ 小→低偏差高方差（过拟合）；$k$ 大→高偏差低方差（欠拟合）；用交叉验证选最优 $k$。
    \item \textbf{距离度量要选对}：欧氏距离最常用；曼哈顿适合高维稀疏；特征缩放是使用任何距离度量的前提。
    \item \textbf{必须标准化}：KNN 基于距离，量纲不一致会导致某个特征完全主导结果，其他特征失效。
    \item \textbf{模型保存}：scaler 和 model 必须配套保存，部署时用同一套预处理参数。
\end{itemize}
\end{document}
